\documentclass[12pt,a4paper]{article}
\usepackage[utf8]{inputenc}
\usepackage[spanish,es-noquoting]{babel}
\usepackage[T1]{fontenc}
\usepackage{lmodern}
\usepackage{geometry}
\geometry{top=2cm, bottom=2cm, left=2.5cm, right=2.5cm}
\usepackage{xcolor}
\usepackage{listings}
\usepackage{enumitem}
\usepackage{fancyhdr}
\usepackage{hyperref}
\usepackage{tikz}
\usetikzlibrary{arrows.meta, positioning, calc, shapes.geometric}

\definecolor{codegreen}{rgb}{0,0.6,0}
\definecolor{backcolour}{rgb}{0.95,0.95,0.92}
\definecolor{lockred}{rgb}{0.8,0,0}

\lstset{
    language=C,
    backgroundcolor=\color{backcolour},
    commentstyle=\color{codegreen},
    keywordstyle=\color{blue},
    basicstyle=\ttfamily\small,
    breaklines=true,
    frame=single,
    numbers=left,
    tabsize=4,
    extendedchars=true,
    showstringspaces=false,
    inputencoding=utf8,
    literate={á}{{\'a}}1 {é}{{\'e}}1 {í}{{\'i}}1 {ó}{{\'o}}1 {ú}{{\'u}}1 {ñ}{{\~n}}1 {¡}{{\textexclamdown}}1 {¿}{{\textquestiondown}}1
}

\pagestyle{fancy}
\fancyhf{}
\lhead{Monitorías Sistemas Operativos 2026-I}
\rhead{Capítulo 4: Sincronización - Mutex}
\cfoot{\thepage}

\title{\textbf{Guía de Aprendizaje: Exclusión Mutua (Mutex)}}
\author{Malak Sanchez \\ \small Monitorías de Sistemas Operativos}
\date{\today}

\begin{document}

\maketitle

\section{El Problema: Condición de Carrera}
Cuando múltiples hilos acceden y modifican un recurso compartido (como una variable global) al mismo tiempo, el resultado final puede depender del orden impredecible de ejecución. Esto se conoce como \textbf{Condición de Carrera} (Race Condition).

\section{La Solución: Mutex}
Un \textbf{Mutex} (MUTual EXclusion) es un mecanismo de sincronización que asegura que solo un hilo a la vez pueda acceder a una sección crítica de código. Es como una llave única: solo el hilo que tiene la llave puede entrar a la habitación; los demás deben esperar afuera.

\begin{center}
\begin{tikzpicture}[
    font=\small,
    >=Stealth,
    thread/.style={
        draw,
        thick,
        rounded corners=2pt,
        fill=blue!10,
        minimum width=3cm,
        minimum height=1.2cm,
        align=center
    },
    resource/.style={
        draw,
        thick,
        rounded corners=2pt,
        fill=green!20,
        minimum width=4cm,
        minimum height=1.5cm,
        align=center
    },
    lock/.style={
        draw,
        thick,
        circle,
        fill=red!55,
        minimum size=2cm,
        align=center
    }
]

    % Mutex
    \node[lock] (mutex) at (0,0) {\textbf{LOCKED}};

    % Threads
    \node[thread] (t1) at (-4,2) {Hilo A\\(En ejecución)};
    \node[thread] (t2) at (4,2) {Hilo B\\(Bloqueado)};

    % Resource
    \node[resource] (res) at (0,-3) {Recurso Compartido\\(Variable Global)};

    % Arrows
    \draw[->, ultra thick, blue]
        (t1.south east) to[bend left=10]
        node[midway, above, sloped] {Posee la llave}
        (mutex.north west);

    \draw[->, ultra thick, gray]
        (t2.south west) to[bend right=10]
        node[midway, above, sloped] {Espera}
        (mutex.north east);

    \draw[->, thick]
        (mutex.south) -- (res.north);

\end{tikzpicture}
\end{center}

\newpage

\section{Funciones Principales}
\begin{lstlisting}[numbers=none]
#include <pthread.h>

pthread_mutex_t mutex;

// Inicializacion
pthread_mutex_init(&mutex, NULL);

// Bloquear (Entrar a la seccion critica)
pthread_mutex_lock(&mutex);

// Desbloquear (Salir de la seccion critica)
pthread_mutex_unlock(&mutex);

// Destruir
pthread_mutex_destroy(&mutex);
\end{lstlisting}

\section{Ejemplo: Contador Seguro}
\begin{lstlisting}
#include <stdio.h>
#include <pthread.h>

int contador = 0;
pthread_mutex_t lock;

void* incrementar(void* arg) {
    for(int i = 0; i < 1000000; i++) {
        pthread_mutex_lock(&lock);
        contador++;
        pthread_mutex_unlock(&lock);
    }
    return NULL;
}

int main() {
    pthread_t t1, t2;
    pthread_mutex_init(&lock, NULL);

    pthread_create(&t1, NULL, incrementar, NULL);
    pthread_create(&t2, NULL, incrementar, NULL);

    pthread_join(t1, NULL);
    pthread_join(t2, NULL);

    printf("Valor final del contador: %d\n", contador);
    pthread_mutex_destroy(&lock);
    return 0;
}
\end{lstlisting}

\section{Ejercicios}
\begin{enumerate}
    \item ¿Qué sucede si olvida llamar a \texttt{pthread\_mutex\_unlock}? Pruebe el comportamiento.
    \item Implemente un programa donde varios hilos retiren dinero de una cuenta bancaria compartida, asegurando que el saldo nunca sea negativo.
\end{enumerate}

\end{document}
